\section*{Chapitre \ref{ch:b2:planet-formation-meteorites} --- Formation des planètes et météorites}

\begin{solution}{exo:b2:planet-formation-meteorites:1}
Gaz~: \qty{2200}{g.cm^{-2}} à \qty{1}{au}~; à \qty{5.2}{au},
$2200\times5.2^{-3/2} = 2200/11.86 = \qty{186}{g.cm^{-2}}$. Solides~:
$33\times0.33 = \qty{10.9}{g.cm^{-2}}$ à \qty{1}{au}~;
$33/11.86 = \qty{2.78}{g.cm^{-2}}$ à \qty{5.2}{au}. À \qty{5.2}{au}, le gaz est
$186/2.78 = 67$ fois plus lourd que les solides.
\end{solution}

\begin{solution}{exo:b2:planet-formation-meteorites:2}
Aluminium (\qty{1653}{K}, ultraréfractaire), magnésium (\qty{1336}{K}, courant),
fer (\qty{1334}{K}, courant), sodium (\qty{958}{K}, modérément volatil), soufre
(\qty{664}{K}, volatil), oxygène sous forme de glace d'eau (\qty{180}{K}, très
volatil). Les CAI, premiers condensats, sont enrichis en aluminium (et en calcium
et titane)~; la Terre est appauvrie en éléments volatils, sodium, soufre et eau,
par rapport au magnésium et à l'aluminium.
\end{solution}

\begin{solution}{exo:b2:planet-formation-meteorites:3}
Chutes~: \omterm{def:b2:planet-formation-meteorites:chondrule}{chondrites} ordinaires $856/1107 = \qty{77.3}{\percent}$, fers
$49/1107 = \qty{4.4}{\percent}$. Trouvailles~: $38647/44534 = \qty{86.8}{\percent}$ et
$1020/44534 = \qty{2.3}{\percent}$. Trouvailles antarctiques~: $22096/44534 =
\qty{49.6}{\percent}$. Les chutes~: toute chute observée est recueillie quelle que soit
sa classe, alors qu'une trouvaille dépend de la façon dont la \omterm{def:b2:planet-formation-meteorites:fall-find}{météorite} survit au sol,
de son aspect et de l'endroit où l'on cherche, et la moitié des trouvailles proviennent
d'un seul type de lieu.
\end{solution}

\begin{solution}{exo:b2:planet-formation-meteorites:4}
$\Delta t = 1.035\ln(\num{5.23e-5}/\num{1.0e-5}) = 1.035\ln 5.23 =
1.035\times1.654 = \qty{1.71}{Ma}$ après les CAI.
\end{solution}

\begin{solution}{exo:b2:planet-formation-meteorites:5}
$4\pi\Sigma_1 r_1^2 = 4\pi\times\num{2.2e4}\times(\num{1.496e11})^2 =
\qty{6.19e27}{kg}$. À l'intérieur de \qty{30}{au}~: $\times\sqrt{30} = 5.48$, soit
\qty{3.39e28}{kg} $= 0.017$ masse solaire~; à l'intérieur de \qty{100}{au}~: $\times10$,
\qty{6.19e28}{kg} $= 0.031$ masse solaire. Les deux valeurs se situent dans l'intervalle observé,
plutôt vers sa borne supérieure~: la nébuleuse était un disque ordinaire.
\end{solution}

\begin{solution}{exo:b2:planet-formation-meteorites:6}
$r = (120/T)^{7/3}$ au~: pour l'eau, $(120/182)^{7/3} = 0.659^{7/3} =
\qty{0.38}{au}$~; pour le méthane, $(120/41)^{7/3} = 2.93^{7/3} = \qty{12.3}{au}$.
Le grain en plein soleil est la couche superficielle du disque~; le plan médian, à l'abri des poussières qui le
surmontent, est plus froid, si bien que la glace subsiste plus près de l'étoile. Les \omterm{def:b2:planet-formation-meteorites:planetesimal}{planétésimaux} se sont formés dans le
plan médian, de sorte que l'estimation du plan médian s'applique, à condition que le disque soit passif~; le chauffage par
accrétion repousse la \omterm{def:b2:planet-formation-meteorites:snow-line}{ligne des glaces} vers l'extérieur.
\end{solution}

\begin{solution}{exo:b2:planet-formation-meteorites:7}
À \qty{1}{au}~: $r^2 = \num{2.238e22}$, $\Sigma_p = \qty{109}{kg.m^{-2}}$, donc
$2\pi b r^2\Sigma_p = 31.42\times\num{2.238e22}\times109 = \num{7.66e25}$~;
$(\num{7.66e25})^{3/2} = \num{6.71e38}$, divisé par
$(3\times\num{1.989e30})^{1/2} = \num{2.443e15}$~: \qty{2.75e23}{kg} $= 0.046$ masse
terrestre. À \qty{5.2}{au}~: $r^2 = \num{6.05e23}$,
$\Sigma_p = \qty{27.8}{kg.m^{-2}}$, $2\pi b r^2\Sigma_p = \num{5.29e26}$,
$M = \num{1.22e40}/\num{2.443e15} = \qty{4.98e24}{kg} = 0.83$ masse terrestre. Avec
\num{3.5} fois plus de solides, $M$ croît comme $3.5^{3/2} = 6.55$~: 5.5 masses
terrestres.
\end{solution}

\begin{solution}{exo:b2:planet-formation-meteorites:8}
Par kilogramme, $\num{1.01e19}$ atomes, donc
$\num{1.01e19}\times\num{5.00e-13} = \qty{5.05e6}{J.kg^{-1}}$. L'élévation de
température serait $\num{5.05e6}/837 = \qty{6030}{K}$. Chauffer de \num{250} à
\qty{1700}{K} demande $837\times1450 = \qty{1.21e6}{J.kg^{-1}}$~: l'aluminium 26 en
fournit environ quatre fois plus.
\end{solution}

\begin{solution}{exo:b2:planet-formation-meteorites:9}
$\num{5.07e10}/\num{7.49e7} = 677$. La durée de vie moyenne du fer 60 vaut
$2.62/\ln2 = \qty{3.78}{Ma}$. Puissances initiales par gramme~:
$\num{5.07e10}/1.035 = \num{4.90e10}$ et $\num{7.49e7}/3.78 = \num{1.98e7}$ erg par gramme et par
million d'années, soit un rapport d'environ \num{2470}~: l'aluminium 26 domine de loin.
\end{solution}

\begin{solution}{exo:b2:planet-formation-meteorites:10}
Les puissances sont égales lorsque $\num{4.90e10}\,e^{-t/1.035} = \num{1.98e7}\,e^{-t/3.78}$, soit
$t(1/1.035 - 1/3.78) = \ln(\num{4.90e10}/\num{1.98e7}) = \ln 2475 = 7.81$~; avec $0.966 - 0.265 = 0.702$ par million d'années,
$t = \qty{11.1}{Ma}$. D'ici là,
il reste au fer 60 $\num{7.49e7}\,e^{-11.1/3.78} = \num{7.49e7}\times0.053 =
\qty{4.0e6}{erg.g^{-1}}$, soit \qty{400}{J.kg^{-1}}~: des milliers de fois trop peu
pour fondre quoi que ce soit.
\end{solution}

\begin{solution}{exo:b2:planet-formation-meteorites:11}
$\tau = 1.035\times\num{3.156e13} = \qty{3.267e13}{s}$ et
$H_0\tau = \qty{6.21e6}{J.kg^{-1}}$. Pour \qty{1200}{K}~:
$t_a^* = 1.035\ln\bigl(\num{6.21e6}/(837\times950)\bigr) = 1.035\ln7.81 =
\qty{2.13}{Ma}$~; pour \qty{1700}{K}~: $1.035\ln\bigl(\num{6.21e6}/(837\times1450)\bigr)
= 1.035\ln5.12 = \qty{1.69}{Ma}$. Diviser $H_0$ par deux retranche $\tau\ln2 =
\qty{0.72}{Ma}$ aux deux, puisque $t_a^* = \tau\ln H_0 + \tau\ln(\tau/c(T_c -
T_0))$, et seul le premier terme change~: le décalage vaut une demi-vie, quelle que soit
la température critique.
\end{solution}

\begin{solution}{exo:b2:planet-formation-meteorites:12}
$\kappa = 2.1/(3300\times837) = \qty{7.60e-7}{m^2.s^{-1}}$ et
$\sqrt{\kappa\tau} = \sqrt{\num{7.60e-7}\times\num{3.267e13}} = \qty{4980}{m}$, soit environ \qty{5}{km}.
Le centre d'un corps de \qty{100}{km} de diamètre est à \qty{50}{km} de profondeur, soit dix fois cette
distance~: sa chaleur ne peut pas s'échapper pendant que l'aluminium 26 se désintègre, et la proposition
s'applique. Un corps de \qty{2}{km} de diamètre perd sa chaleur aussi vite qu'il la produit. La fusion
exige un rayon de quelques fois \qty{5}{km}, soit un diamètre d'environ \qtyrange{10}{20}{km} ou plus,
comme le trouvent les modèles.
\end{solution}

\begin{solution}{pb:b2:planet-formation-meteorites:1}
\textbf{1.} Une chute~: son arrivée a été vue, et elle a été ramassée le lendemain matin. Elle n'est pas altérée, et elle s'ajoute à
l'échantillon non biaisé de ce qui atteint la Terre.

\textbf{2.} $856/1107 = 0.77$~: environ trois chances sur quatre.

\textbf{3.} \omterm{def:b2:planet-formation-meteorites:chondrule}{Chondres} de \qty{0.3}{mm} de diamètre, \qty{9}{\percent} de métal et olivine à \qty{18}{\percent} de fayalite~: une
\omterm{def:b2:planet-formation-meteorites:chondrule}{chondrite} H (\qty{10}{\percent} de métal, Fa 16--20).

\textbf{4.} Des \omterm{def:b2:planet-formation-meteorites:chondrule}{chondres} estompés et des minéraux de composition uniforme signalent un
métamorphisme thermique~: types 5 à 6. Les \omterm{def:b2:planet-formation-meteorites:chondrule}{chondrites} ordinaires ont atteint au plus environ
\qty{1220}{K}, soit \qty{947}{\degreeCelsius}.

\textbf{5.} Les \omterm{def:b2:planet-formation-meteorites:chondrule}{chondrites} ordinaires ne contiennent que \qtyrange{0.1}{1}{\percent} de CAI en volume. La
pierre est faite de \omterm{def:b2:planet-formation-meteorites:chondrule}{chondres} et de matrice assemblés après les CAI, qui sont les solides les plus anciens~:
toute roche qui les contient est plus jeune.

\textbf{6.} Chaque minéral a reçu à sa formation le même magnésium, plus de
l'aluminium 26 en proportion de son aluminium~; l'aluminium 26 s'est depuis
transformé en magnésium 26, si bien que plus un minéral avait d'aluminium par
magnésium, plus il a gagné de magnésium 26.

\textbf{7.} $(0.139590 - 0.139500)/(15.0 - 0.05) = \num{0.000090}/14.95 =
\num{6.02e-6}$~: le rapport $\ce{^{26}Al}/\ce{^{27}Al}$ du \omterm{def:b2:planet-formation-meteorites:chondrule}{chondre} lorsqu'il
s'est formé.

\textbf{8.} $0.139500 + \num{6.02e-6}\times3.00 = 0.139518$~: sur la droite.

\textbf{9.} $\Delta t = 1.035\ln(\num{5.23e-5}/\num{6.02e-6}) = 1.035\ln8.69 =
1.035\times2.162 = \qty{2.24}{Ma}$.

\textbf{10.} $4567.5 - 2.2 = \qty{4565.3}{\Ma}$.

\textbf{11.} $\num{5.23e-5}\,e^{-10/1.035} = \num{5.23e-5}\times\num{6.4e-5} =
\num{3.3e-9}$, bien trop peu de magnésium 26 pour être mesuré. L'hafnium 182 (demi-vie
\qty{8.9}{Ma}) ou le manganèse 53 (\qty{3.7}{Ma}) conviennent aux événements survenus environ dix millions d'années
après les CAI.

\textbf{12.} $\kappa = 2.1/(3300\times837) = \qty{7.60e-7}{m^2.s^{-1}}$~;
$R^2/\kappa = \num{4e8}/\num{7.60e-7} = \qty{5.26e14}{s} = \qty{16.7}{Ma}$, seize fois la durée de vie moyenne de
l'aluminium 26~: le centre conserve sa chaleur.

\textbf{13.} $\num{1.9e-7}\,e^{-2.4/1.035} = \num{1.9e-7}\times0.0984 =
\qty{1.87e-8}{W.kg^{-1}}$.

\textbf{14.} $\num{1.87e-8}\times\num{3.267e13} = \qty{6.11e5}{J.kg^{-1}}$.

\textbf{15.} $\Delta T = \num{6.11e5}/837 = \qty{730}{K}$~; maximum
$250 + 730 = \qty{980}{K}$.

\textbf{16.} Au-dessus de \qty{750}{K}, la roche a donc été frittée et métamorphisée~; en dessous de \qty{1200}{K}, ni le
métal ni le sulfure n'ont fondu. La pierre est une \omterm{def:b2:planet-formation-meteorites:chondrule}{chondrite} chauffée, d'un \omterm{def:b2:planet-formation-meteorites:type}{type pétrologique} élevé, en dessous des
\qty{1220}{K} de la question~4~: cela est cohérent.

\textbf{17.} $1.035\ln\bigl(\num{6.21e6}/(837\times950)\bigr) = \qty{2.13}{Ma}$~: environ
\qty{0.3}{Ma} plus tôt.

\textbf{18.} L'élévation est $\Delta T(1 - e^{-(t - t_a)/\tau})$~; \qty{90}{\percent}
exige $e^{-(t - t_a)/\tau} = 0.1$, $t - t_a = 1.035\ln10 = \qty{2.38}{Ma}$, soit
\qty{4.8}{Ma} après les CAI.

\textbf{19.} $\Delta T = 7420\,e^{-0.8/1.035} = 7420\times0.462 = \qty{3430}{K}$, maximum
\qty{3680}{K}, bien au-dessus des \qty{1700}{K} auxquels les silicates sont fondus. Le corps a
fondu (la chaleur de fusion a maintenu sa température plus basse) et s'est différencié en un noyau
métallique, un manteau silicaté et une croûte.

\textbf{20.} $e^{(2.4 - 0.8)/1.035} = e^{1.546} = 4.7$ fois plus.

\textbf{21.} Le zinc et le plomb se condensent vers \qty{727}{K}, éléments modérément volatils~; le
magnésium (\qty{1336}{K}) et l'aluminium (\qty{1653}{K}) se condensent bien plus tôt. Les solides
de la région de la pierre se sont rassemblés avant que le gaz ne soit assez refroidi pour que tout
le zinc et le plomb se condensent.

\textbf{22.} $33\times0.33\times2.5^{-3/2} = 10.89/3.95 = \qty{2.76}{g.cm^{-2}}$,
\qty{27.6}{kg.m^{-2}}.

\textbf{23.} $r^2 = (3.74\times10^{11})^2 = \num{1.40e23}$~;
$2\pi\times5\times\num{1.40e23}\times27.6 = \num{1.21e26}$~;
$(\num{1.21e26})^{3/2}/\num{2.443e15} = \qty{5.45e23}{kg}$, soit $0.091$ masse terrestre. Le corps
parent, $\tfrac43\pi(\num{2e4})^3\times3300 = \qty{1.1e17}{kg}$, est cinq millions de fois plus
léger~: un \omterm{def:b2:planet-formation-meteorites:planetesimal}{planétésimal}, et non un embryon.

\textbf{24.} Le centre du corps parent a atteint une température maximale d'environ
\textbf{\qty[text-series-to-math=true]{707}{\degreeCelsius}} ($980 - 273$)~: assez chaud pour recristalliser la
\omterm{def:b2:planet-formation-meteorites:chondrule}{chondrite}, trop froid pour fondre son métal.
\end{solution}
